\begin{table}[H]
\centering
\begin{threeparttable}
\caption{Textual Balance of the Final Reading Materials}
\label{tab:text_balance}
\footnotesize
\begin{tabular}{lccccc}
\toprule
Group & $N$ & Word count & Words/sentence & Chars/word & Flesch-Kincaid\\
\midrule
\addlinespace[0.3em]
\multicolumn{6}{l}{\textit{Panel A. Article pools}}\\
Pro-China & 24 & 722 & 25.19 & 5.01 & 14.10\\
Anti-China & 24 & 898 & 24.58 & 5.00 & 13.79\\
Apolitical China & 17 & 846 & 25.75 & 4.88 & 13.60\\
Non-China control & 24 & 989 & 24.95 & 4.91 & 13.59\\
ANOVA $p$-value & -- & 0.001 & 0.925 & 0.084 & 0.895\\
\addlinespace[0.45em]
\multicolumn{6}{l}{\textit{Panel B. Average assigned 24-reading bundle by group}}\\
Pro low & 1061 & 924 & 25.01 & 4.94 & 13.72\\
Pro high & 1061 & 856 & 25.13 & 4.96 & 13.87\\
Anti low & 1061 & 967 & 24.89 & 4.94 & 13.65\\
Anti high & 1061 & 945 & 24.80 & 4.96 & 13.71\\
Apolitical China & 1061 & 917 & 25.35 & 4.90 & 13.60\\
Control & 1060 & 989 & 24.95 & 4.91 & 13.59\\
ANOVA $p$-value & -- & <0.001 & <0.001 & <0.001 & <0.001\\
\bottomrule
\end{tabular}
\begin{tablenotes}
\footnotesize
\item Note: Panel A reports article-level averages by content bank. Panel B reports participant-level averages of the 24 assigned readings using the implemented schedule design and the realized randomized sample. Higher Flesch-Kincaid values indicate more difficult reading material. These diagnostics assess non-substantive textual comparability rather than political valence.
\end{tablenotes}
\end{threeparttable}
\end{table}
